\documentclass[12pt]{article}
\usepackage[utf8]{inputenc}
\usepackage[T1]{fontenc}
\usepackage{hyperref}
\title{Mandandanyi Language Resources for K-12 Curriculum Development}
\author{Compiled by Claude Sonnet 5.5 (claude-sonnet-5-5) for Andr\'as Kornai}
\date{}
\begin{document}
\maketitle

\section{Introduction}
Mandandanyi (also spelled Mandandanji; ISO 639-3 zmk, Glottocode mand1417) is a Pama-Nyungan language of southern Queensland, in the Maranoa and Balonne river region. It is extinct or critically endangered. AIATSIS notes that Mandandanyi, Gunggari and Kogai may refer to the same language, and sources disagree on this.

\section{Primary Educational Resources}
No primers, textbooks or school programmes for Mandandanyi were found in search results.

\section{Linguistic and Reference Materials}
\begin{itemize}
\item Wikipedia on the Mandandanji people: \url{https://en.wikipedia.org/wiki/Mandandanji}
\item Endangered Languages Project bibliography: \url{https://endangeredlanguages.com/lang/3136/bibliography}
\item AIATSIS Austlang entry for the related Gunggari (D37): \url{https://collection.aiatsis.gov.au/austlang/language/d37}
\item AIATSIS collection entry: \url{https://collection.aiatsis.gov.au/node/1508013}
\item University of Queensland manuscripts entry, Mandandanji people D44: \url{https://manuscripts.library.uq.edu.au/index.php/mandandanji-people-d44}
\item Bidjara language (Maric group context): \url{https://en.wikipedia.org/wiki/Bidjara_language}
\end{itemize}

\section{Religious and Community Resources}
\begin{itemize}
\item Mandandanji community website (traditional custodians): \url{https://www.mandandanji.com.au/about-us/}
\item Handbook entry on the Mandandanji: \url{https://hgcosmos.org/assets/culture/AZ/M/Mandandanji.pdf}
\end{itemize}

\section{Digital and Audio Resources}
No audio or digital language resources were found in search results.

\section{Curriculum Development Context}
Little linguistic documentation appears to be available. A revitalization curriculum would probably begin with cultural and place-name content, led by the Mandandanji community. Archival wordlists, including those listed under Gunggari and Kogai, would need community consultation before use.
\end{document}